% GENERATED FILE. Edit canonical lessons, then run npm run generate:books.
% canonical-source-hash: fnv1a64:8595b21f5d068867
% canonical-lessons: GE-C314-zaubern, GE-C314-verzaubern, GE-C314-verwandeln, GE-C314-fluch, GE-C314-drache

\chapter{Verzaubert}
\label{ch:ge-verzaubert}
\hlchaptermodality{314}

\begin{chapteropening}
\textbf{By the end of this chapter:} \emph{I can tell what happens by magic: to conjure, to bewitch, to turn into, a curse and a dragon.}

Complete the last lesson of chapter 314: I can tell what happens by magic: to conjure, to bewitch, to turn into, a curse and a dragon.
\end{chapteropening}

\section[zaubern]{zaubern — to conjure, to do magic}

\label{lesson:GE-C314-zaubern}



\begin{quote}
Before the new one: say the German for the witch, then the German for the wizard.
\end{quote}

\subsection*{You'll want to know: zaubern}
\textbf{zaubern} — \textquotedblleft{}to conjure, to do magic\textquotedblright{}.

\begin{quote}
Der Zauberer zauberte eine Blume aus dem Hut.
\end{quote}

\begin{quote}
\emph{The magician conjured a flower out of the hat.}
\end{quote}

Regular: \emph{zauberte}, \emph{hat gezaubert}.

\begin{cousinweb}
\emph{zaubern} is inherited from Middle High German \emph{zoubern} (EtymDB), built on \emph{Zauber} \textquotedblleft{}magic\textquotedblright{}, Proto-Germanic \emph{taubrą}.
\end{cousinweb}

\begin{grammarlens}[title={What you've built}]
One more word for telling, explaining and describing at B1.
\end{grammarlens}

\subsection*{Your turn}
\begin{itemize}
\raggedright
  \item \emph{Say it:} \emph{zaubern}
  \item \emph{Say it:} \emph{zaubern}, once more
  \item \emph{Say it:} a sentence of your own with \emph{zaubern}
  \item \emph{Recall:} say the German for the witch, then the German for the wizard, then say \emph{zaubern} again
\end{itemize}

\begin{tcolorbox}[breakable,colback=teal!4,colframe=teal!35!black,title={Before you move on}]
What does zaubern mean? (\textquotedblleft{}To conjure, to do magic\textquotedblright{}.) Say it once more.
\end{tcolorbox}
\section[verzaubern]{verzaubern — to bewitch, to put a spell on; to charm}

\label{lesson:GE-C314-verzaubern}



\begin{quote}
Before the new one: say the German for the wizard, then the German for to conjure.
\end{quote}

\subsection*{You'll want to know: verzaubern}
\textbf{verzaubern} — \textquotedblleft{}to bewitch, to put a spell on; to charm\textquotedblright{}.

\begin{quote}
Die Hexe verzauberte das Mädchen.
\end{quote}

\begin{quote}
\emph{The witch put a spell on the girl.}
\end{quote}

Also figurative: \emph{Die Musik verzauberte alle}.

\begin{cousinweb}
\emph{verzaubern} is built on \emph{zaubern}, Middle High German \emph{zoubern} (EtymDB), with \emph{ver-}: to change by magic.
\end{cousinweb}

\begin{grammarlens}[title={What you've built}]
One more word for telling, explaining and describing at B1.
\end{grammarlens}

\subsection*{Your turn}
\begin{itemize}
\raggedright
  \item \emph{Say it:} \emph{verzaubern}
  \item \emph{Say it:} \emph{verzaubern}, once more
  \item \emph{Say it:} a sentence of your own with \emph{verzaubern}
  \item \emph{Recall:} say the German for the wizard, then the German for to conjure, then say \emph{verzaubern} again
\end{itemize}

\begin{tcolorbox}[breakable,colback=teal!4,colframe=teal!35!black,title={Before you move on}]
What does verzaubern mean? (\textquotedblleft{}To bewitch, to put a spell on; to charm\textquotedblright{}.) Say it once more.
\end{tcolorbox}
\section[verwandeln]{verwandeln — to turn into, to transform}

\label{lesson:GE-C314-verwandeln}



\begin{quote}
Before the new one: say the German for to conjure, then the German for to bewitch.
\end{quote}

\subsection*{You'll want to know: verwandeln}
\textbf{verwandeln} — \textquotedblleft{}to turn into, to transform\textquotedblright{}.

\begin{quote}
Die Hexe verwandelte ihn in einen Frosch.
\end{quote}

\begin{quote}
\emph{The witch turned him into a frog.}
\end{quote}

\emph{In etwas verwandeln}; reflexive, \emph{sich verwandeln} — to change into.

\begin{cousinweb}
\emph{verwandeln} is inherited from Middle High German \emph{verwandelen} (EtymDB), built on \emph{wandeln} \textquotedblleft{}to change\textquotedblright{}, a relative of \emph{wenden} \textquotedblleft{}to turn\textquotedblright{}.
\end{cousinweb}

\begin{grammarlens}[title={What you've built}]
One more word for telling, explaining and describing at B1.
\end{grammarlens}

\subsection*{Your turn}
\begin{itemize}
\raggedright
  \item \emph{Say it:} \emph{verwandeln}
  \item \emph{Say it:} \emph{verwandeln}, once more
  \item \emph{Say it:} a sentence of your own with \emph{verwandeln}
  \item \emph{Recall:} say the German for to conjure, then the German for to bewitch, then say \emph{verwandeln} again
\end{itemize}

\begin{tcolorbox}[breakable,colback=teal!4,colframe=teal!35!black,title={Before you move on}]
What does verwandeln mean? (\textquotedblleft{}To turn into, to transform\textquotedblright{}.) Say it once more.
\end{tcolorbox}
\section[der Fluch]{der Fluch — the curse}

\label{lesson:GE-C314-fluch}



\begin{quote}
Before the new one: say the German for to bewitch, then the German for to turn into.
\end{quote}

\subsection*{You'll want to know: der Fluch}
\textbf{der Fluch} — \textquotedblleft{}the curse\textquotedblright{}.

\begin{quote}
Ein Fluch lag auf der alten Burg.
\end{quote}

\begin{quote}
\emph{A curse lay on the old castle.}
\end{quote}

Plural: \emph{die Flüche}. Also a swear word.

\begin{cousinweb}
\emph{Fluch} is inherited from Old High German \emph{fluoh} (EtymDB), related to the verb \emph{fluchen} \textquotedblleft{}to curse, to swear\textquotedblright{}.
\end{cousinweb}

\begin{grammarlens}[title={What you've built}]
One more word for telling, explaining and describing at B1.
\end{grammarlens}

\subsection*{Your turn}
\begin{itemize}
\raggedright
  \item \emph{Say it:} \emph{der Fluch}
  \item \emph{Say it:} \emph{der Fluch}, once more
  \item \emph{Say it:} a sentence of your own with \emph{der Fluch}
  \item \emph{Recall:} say the German for to bewitch, then the German for to turn into, then say \emph{der Fluch} again
\end{itemize}

\begin{tcolorbox}[breakable,colback=teal!4,colframe=teal!35!black,title={Before you move on}]
What does der Fluch mean? (\textquotedblleft{}The curse\textquotedblright{}.) Say it once more.
\end{tcolorbox}
\section[der Drache]{der Drache — the dragon}

\label{lesson:GE-C314-drache}



\begin{quote}
Before the new one: say the German for to turn into, then the German for the curse.
\end{quote}

\subsection*{You'll want to know: der Drache}
\textbf{der Drache} — \textquotedblleft{}the dragon\textquotedblright{}.

\begin{quote}
Der Drache wohnte in einem Berg.
\end{quote}

\begin{quote}
\emph{The dragon lived in a mountain.}
\end{quote}

Weak noun: \emph{den Drachen}. \emph{Der Drachen}, with \emph{-n}, is a kite.

\begin{cousinweb}
\emph{Drache} comes from Old High German \emph{trahho} (EtymDB), an early loan from Latin \emph{dracō} \textquotedblleft{}a serpent, a dragon\textquotedblright{} (Lewis \& Short), itself from Greek \emph{drakōn}.
\end{cousinweb}

\begin{grammarlens}[title={What you've built}]
One more word for telling, explaining and describing at B1.
\end{grammarlens}

\subsection*{Your turn}
\begin{itemize}
\raggedright
  \item \emph{Say it:} \emph{der Drache}
  \item \emph{Say it:} \emph{der Drache}, once more
  \item \emph{Say it:} a sentence of your own with \emph{der Drache}
  \item \emph{Recall:} say the German for to turn into, then the German for the curse, then say \emph{der Drache} again
\end{itemize}

\begin{tcolorbox}[breakable,colback=teal!4,colframe=teal!35!black,title={Before you move on}]
What does der Drache mean? (\textquotedblleft{}The dragon\textquotedblright{}.) Say it once more.
\end{tcolorbox}
